\section{Finite bounds for widths two and four}\label{app:wide-mask}

The measured wider encoders use the same Newton precode $B_8$ and two
permutations, without a final unpermuted scan. We extend the local rank
bound by recording the support mask and rank of the preceding $w$ outputs.
Fix the shortened source space $V_S$ and quotient by the zero equations
imposed so far. Let $M\subseteq\{1,\ldots,w\}$ record the positions of the
required nonzero feedback forms, and let $R$ be their span dimension in
this quotient. The states are $(\varnothing,0)$ and
$(M,R)$ with $1\le R\le|M|$. There are $1+w2^{w-1}$ possible states.
Choose positive continuation values ordered by rank:
$v_{M,R+1}\ge v_{M,R}$. This is a condition on the witness, not a claim
that the exact continuation count increases with rank.

Drop the oldest state and shift the remaining mask to obtain $D$, with
its new-coordinate bit zero. Let $m=|D|$. With zero input, the output
form lies in the span of the state forms. A \emph{contraction} is a
zero output whose form is a nonzero element of this span: the new zero
equation lowers the rank of the state span by one. The following
successor envelopes bound arbitrary linear dependencies:
\begin{center}
\begin{tabular}{lcc}
\toprule
Event & Successor mask & Rank bound\\
\midrule
Zero output form ($Z$) & $D$ & $\min(R,m)$\\
Nonzero output in state span ($I$) & $D\cup\{\mathrm{new}\}$ & $\min(R,m+1)$\\
Output outside state span ($O$) & $D\cup\{\mathrm{new}\}$ & $\min(R,m)+1$\\
Contract a state-span form ($C$) & $D$ & $\min(R-1,m)$\\
\bottomrule
\end{tabular}
\end{center}
The last event is allowed only for $R\ge2$ and a surviving contraction,
one after which every required form remains nonzero. If the new zero
equation makes a required form zero, that branch is discarded. A zero
output whose form lies outside the state span preserves the old state
rank and is bounded by the $Z$ envelope. Rank ordering lets each
envelope use its displayed upper rank. Write $v_Z,v_I,v_O,v_C$ for the
continuation values at the successor states of the four events.

At $R>0$, fresh independent coefficients in $\{1,\ldots,q\}$ make
the output form zero with probability at most $\pi=q^{-R}$.
Condition on all but $R$ pivot coefficients: at most one pivot assignment
cancels the form. With zero input and $R=|M|$, cancellation is impossible
because the state forms are independent and every coefficient is nonzero.
Define
\[
 \Theta_0=yv_I+\mathbf1_{R\ge2}\tau^{-1}v_C,\qquad
 \Theta_{\rm ind}=yv_O+\tau^{-1}v_Z,\qquad
 \Theta_{\rm dep}=(1-\pi)\Theta_0+\pi v_Z.
\]
At rank one, the state span is spanned by a single required form, so a
contraction would make that form zero; this is why $\Theta_0$ has no
zero branch when $R=1$. The $O$ and $I$ successors share a mask and $O$
has at least the rank of $I$, so rank ordering gives $v_O\ge v_I$;
likewise $v_Z\ge v_C$. Hence $\Theta_{\rm ind}\ge\Theta_0$.
Maximizing over the unknown cancellation probability in $[0,\pi]$ gives,
in analogy with~\eqref{eq:width-one-quotient},
\[
 T_{M,R}(v)=\begin{cases}
 \Theta_0+x\max(\Theta_{\rm ind},\Theta_{\rm dep}),&R=|M|,\\
 \max(\Theta_0,\Theta_{\rm dep})+x\max(\Theta_{\rm ind},\Theta_{\rm dep}),&R<|M|,
 \end{cases}
\]
and $T_0(v)=v_0+xyv_{\{\mathrm{new}\},1}$.
After each step, take the maximum over smaller ranks at a fixed mask.
This closure preserves the ordering. If $T(v)\le\lambda v$, its
closure satisfies the same inequality. With $v_0=1$, put
$\Gamma(a)=\max_i[T^a(\mathbf1)]_i/v_i$. Positive homogeneity gives
$[T^n(\mathbf1)]_0\le\Gamma(a)\lambda^{n-a}$, as
in~\eqref{eq:scan-terminal-sixteen} for $a=16$.

All successor bounds hold for arbitrary incoming forms and earlier
nonzero requirements. Lemma~\ref{lem:scan-composition} therefore applies
with these scan bounds and the same $F_s(\tau)/\tau$ source factor.
The inverse scan gives
$h\le(w+1)e$ and $s\le(w+1)h$, capped by $n$.
For $s<32$, compute the source coefficients exactly, use $a=16$, and
cover all feasible intermediate weights by consecutive integer intervals.
For $s\ge32$, use $F_s(\tau)\le Q_8(z,\tau)^{n/16}/z^s$ and $a=0$.
The rational binomial prefactors and convex vertex bounds of
Appendix~\ref{app:all-one} then apply with the wider support constraint.
The dense domain is partitioned recursively into rectangles; each cell
is a leaf, and each leaf is charged its full number of integer profiles.
The leaves cover the entire feasible domain; summing the interval and
leaf bounds with upward rounding gives the complete failure bound.

For $w=2$ and $q=2048$, at forbidden weight $e=829633$ the sparse range
uses 185 intervals and the dense cover has 216 leaves; the total failure
probability has base-two logarithm at most $-57.875553147005$. For $w=4$
and $q=65536$, at $e=901502$ the sparse range uses 180 intervals and
the dense cover has 229 leaves; the total has base-two logarithm at most
$-61.763321686333$. These certificates yield the two rows of
Table~\ref{tab:wide-measured}. Failure exponents there are truncated
toward zero, so the stated bounds are conservative.
